\documentclass[letter,10pt]{book}
%-----------------------------------------------------------
\usepackage[top=0.25in, bottom=0.25in, left=0.55in, right=0.55in]{geometry}
\usepackage{graphicx}
\usepackage{url}
\usepackage{palatino}
\usepackage{pdfrender}
\usepackage{tabularx}
\usepackage[colorlinks=true,linkcolor=black,anchorcolor=black,citecolor=black,filecolor=black,menucolor=black,runcolor=black,urlcolor=black]{hyperref}
\usepackage{wrapfig}
\fontfamily{SansSerif}
\selectfont

\usepackage[T1]{fontenc}
\usepackage
%[ansinew]
[utf8]
{inputenc}

\usepackage{color}
\definecolor{mygrey}{gray}{0.75}
\textheight=9.75in
\raggedbottom

\setlength{\tabcolsep}{0in}
\newcommand{\isep}{-2 pt}
\newcommand{\lsep}{-0.5cm}
\newcommand{\psep}{-0.6cm}
\renewcommand{\labelitemii}{$\circ$}

\pagestyle{empty}
%-----------------------------------------------------------
%Custom commands
\newcommand{\resitem}[1]{\item #1 \vspace{-2pt}}
\newcommand{\resheading}[1]{{\small \colorbox{mygrey}{\begin{minipage}{0.975\textwidth}{\textbf{#1 \vphantom{p\^{E}}}}\end{minipage}}}}
\newcommand{\ressubheading}[3]{
\begin{tabular*}{6.62in}{l @{\extracolsep{\fill}} r}
	\textsc{{\textbf{#1}}} & \textsc{\textit{[#2]}} \\
\end{tabular*}\vspace{-8pt}}
%-----------------------------------------------------------
\begin{document}
\begin{center}
\textsc{\huge Aditya Rastogi}
\end{center}

\begin{minipage}{0.30\textwidth}
{\raggedleft
\href{https://github.com/thunderInfy/}{GitHub Profile} \\
\href{http://cse.iitkgp.ac.in/~arastogi}{Personal Website CSE - IITKGP}\\
}
\end{minipage}%
\begin{minipage}{0.33\textwidth}
\begin{center}
	\includegraphics[scale=0.1]{logo.jpg}
\end{center}
\end{minipage}%
\begin{minipage}{0.33\textwidth}
{\raggedright
Email-id : \textbf{r.aditya0824@gmail.com} \\
Mobile No.: \textbf{917477725222} \\
}
\end{minipage}%
\hspace{0.5cm}\\
\\
\resheading{\textbf{ACADEMIC DETAILS} }\\[\lsep]
\\
\begin{center}
\renewcommand{\arraystretch}{1.5}
\indent \begin{tabular}{ |@{\hskip 0.125in}l @{\hskip 0.125in} |@{\hskip 0.125in}l @{\hskip 0.125in} |@{\hskip 0.125in}l @{\hskip 0.125in} |@{\hskip 0.125in}l @{\hskip 0.25in} |l }
\hline
\textbf{Education} & \textbf{Institute} & \textbf{Year} & \textbf{CGPA / \%} \\
\hline
B. Tech and M. Tech (Dual Degree): &&&\\[-0.5em]
{Computer Science and Engineering} & IIT Kharagpur  & 2016 - 2021 (Expected) & \textbf{9.42 / 10} \\
\hline
\end{tabular}
\end{center}
\hfill 
\hspace{0.5cm}\\
\resheading{\textbf{INTERNSHIPS} }\\[\lsep]
\\\\
\large \textbf{Pattern Matching in Trillion Edge Graphs} \normalsize \\
\emph{Mitacs Globalink Research Intern, University of British Columbia \hfill May'19-Jul'19 } \\ 
\emph{(Advisor: Prof. Matei Ripeanu)} \\[\lsep]
\begin{itemize}
\item Worked in a team to solve the problem of \textbf{approximate pattern matching in graphs} in a distributed systems setting.\\[\psep]
\item  Compared our pattern matching algorithm with TriAD, a distributed RDF engine.
\\[\lsep]
\end{itemize}
\emph{Research Areas:}  \textbf{Algorithms}, \textbf{Distributed Systems},
\textbf{Graph Theory}
\\\\
\large \textbf{Facial Landmarks Detection} \normalsize \\
\emph{Visiting Student Researcher, University of Sydney \hfill Dec'18-Jan'19} \\ 
\emph{(Advisor: Dr. Mehar Khatkar)} \\[\lsep]
\begin{itemize}
	\item  Designed and compared different \textbf{CNN} architectures with target point coordinates for detecting facial landmarks in different fish species.
	\\[\lsep]
\end{itemize}
\emph{Research Areas:}  \textbf{Deep Learning}, \textbf{Landmark Detection},
\textbf{Data Augmentation}
\\\\
\large \textbf{Sensor Diagnostics} \normalsize \\
\emph{Summer Intern, IIT-Kharagpur - Sponsored by Shell India Pvt. Ltd. \hfill May'18-July'18} \\ 
\emph{(Advisor: Prof. Swanand Khare)} \\[\lsep]
\begin{itemize}
\item  Developed a \textbf{Python application} to \textbf{sample from an arbitrary finite-ranged probability density} in one dimension.\\[\psep]
\item Reduced dimensionality of multiple sensors time-series data using \textbf{PCA} and \textbf{autoencoders}.\\[\psep]
\item Worked on \textbf{gaussian-mixture models} and error distributions in general.\\[\lsep]
\end{itemize}
\emph{Research Areas:}  \textbf{Machine Learning}, \textbf{Dimensionality Reduction},  \textbf{Predictive}  \textbf{Analysis}
\\
\hspace{0cm}\\[-0.5cm]\\
\resheading{\textbf{ PROJECTS } }\\[\lsep]
\\\\ \large \href{http://cse.iitkgp.ac.in/~arastogi}{\textbf{Smart Rockets, Algorithmic Visualisations and Web-game Development}} \normalsize\\
\emph{p5-js, HTML, CSS} \\[\lsep]
\begin{itemize}
\item Developed Smart Rockets, a demo of a \textbf{genetic algorithm}. Developed algorithmic visualisations on the web for algorithms like \textbf{Convex Hull}, Voronoi Diagram, Longest Increasing Subsequence etc.
\\[\psep]
\item  Used p5.js for creating \textbf{fractal} designs like trees, koch curves and flow field visualisation of \textbf{perlin noise}.\\[\psep]
\item  Developed casual, fast-paced and puzzle 2D web-games.\\
\end{itemize}
\large \textbf{Visualising Deep Neural Network Architectures - B.Tech Project} \normalsize\\
\emph{p5-js, TensorFlow.js \hfill Aug'19 - Present} \\ 
\emph{(Advisors: Prof. Partha Pratim Chakrabarti and Prof. Aritra Hazra)} \\[\lsep]
\begin{itemize}
	\item Developing a visualisation software for deep neural network architectures.\\[\psep]
	\item Developing collaboratorial and adversarial agents to conduct unsupervised and statistical analysis on top of the neural network architecture. 
	\\[\lsep]
\end{itemize}

\hspace{0.5cm}\\[-0.2cm]
\resheading{\textbf{BLOG POSTS} }\\[\lsep]
\\[-0.5em]
\begin{itemize}
	\item
	\href{https://towardsdatascience.com/elucidating-policy-iteration-in-reinforcement-learning-jacks-car-rental-problem-d41b34c8aec7}{Elucidating \textbf{Policy Iteration} in \textbf{Reinforcement Learning} — Jack’s Car Rental Problem} : \textit{Towards Data Science}
\end{itemize}

\pagebreak

\hspace{0.5cm}\\[-0.2cm]
\resheading{\textbf{CERTIFICATIONS} }\\[\lsep]
\\[-0.5em]
\begin{itemize}
	\item \href{https://www.coursera.org/account/accomplishments/certificate/7BQGV7ZGYBHV}{\textbf{Neural Networks and Deep Learning}} : \textit{Deeplearning.ai - Coursera}\\[\psep]
	\item \href{https://www.coursera.org/account/accomplishments/certificate/Q9UDAMNJ4WJW}{\textbf{Improving Deep Neural Networks: Hyperparameter tuning, Regularization and Optimization}} : \textit{Deeplearning.ai - Coursera}\\[\psep]
	\item \href{https://www.coursera.org/account/accomplishments/certificate/XHD5MM85DNN8}{\textbf{Structuring Machine Learning Projects}} : \textit{Deeplearning.ai - Coursera}\\[\psep]
	\item \href{https://www.coursera.org/account/accomplishments/certificate/WJ6TRA83A89V}{\textbf{Mathematics for Machine Learning: PCA}} : \textit{Imperial College London - Coursera}
\end{itemize}


\hspace{0.5cm}\\[-0.2cm]
\resheading{\textbf{TERM PROJECTS} }\\[\lsep]
\\[-1em]
\begin{itemize}
	\item \href{https://github.com/thunderInfy/Spam-Classifier}{Built a deep neural network in python from scratch to classify messages as spam or ham.}\\[\psep]
	\item \href{https://github.com/thunderInfy/Operating-Systems-Lab-IIT-KGP}{Implemented a memory-resident file system in Operating Systems Lab. Simulated Virtual Memory using Demand Paging, compared CPU scheduling algorithms and implemented a command-line interpreter running on Linux.} \\[\psep]
	\item \href{https://github.com/thunderInfy/Networks-Lab-IIT-KGP}{Developed an API to do reliable communication over an unreliable link using timeouts and acknowledgements. Implemented a simplified FTP version in C.}\\[-2em]
	\item Implemented a single-cycle processor (RISC architecture) in Verilog.\\[\psep]
	\item \href{https://github.com/thunderInfy/compiler-for-tiny-C}{Designed a compiler for tiny-C.}\\[\psep]
	\item \href{https://github.com/thunderInfy/Personal-Library-System-software}{Developed a Java-based software to manage a personal library system.}\\[\psep]
	\item Implemented cryptographic algorithms like \href{https://github.com/thunderInfy/AES-one-round}{AES} and \href{https://github.com/thunderInfy/DES-on-image}{DES} in C from scratch and ran them on images and text.
\end{itemize}
\hspace{0.5cm}\\[-0.2cm]
\resheading{\textbf{SKILLS AND EXPERTISE} }\\[\lsep]
\\[-1em]
\begin{itemize}
\item \textbf{Languages :} Python, C, C++, Java, HTML, CSS\\[\psep]
\item \textbf{Libraries and Tools :} TensorFlow.js, Keras, p5.js, Processing, Numpy, Scipy, Matplotlib
\end{itemize}

\hspace{0.5cm}\\[-0.2cm]
\resheading{\textbf{COURSEWORK INFORMATION} }\\[\lsep]
\\[-1em]
\begin{itemize}
\item \textbf{Completed : }*Programming and Data Structures | *Algorithms - I | Discrete Structures | *Signals and Networks | Machine Learning | Formal Language and Automata Theory | *Switching Circuits and Logic Design | *Software Engineering | Probability and Statistics | *Computer Organization and Architecture | *Compilers | Algorithms - II | Data Analytics | Cryptography and Network Security | *Operating Systems | *Computer Networks | Computational Complexity | Combinatorics and Computing | Quantum Computing - Project Seminar 
\\[\psep]
\item \textbf{Will be completed by May 2020 :}  Theory of Computation | Object-Oriented Systems | Artificial Intelligence | Reinforcement Learning | Image Processing | High Performance in Computer Architecture | Linear Algebra\\[\psep]
\item \textbf{Micro Credit Course :} Computational Intelligence in Cyber Security\\
{*} marked courses have a laboratory component as well.
\end{itemize}

\hspace{0.5cm}\\[-0.2cm]
\resheading{\textbf{POSITIONS OF RESPONSIBILITY} }\\[\lsep]
\\[-1em]
\begin{itemize}
\item \large \textbf{Student Mentor} \hfill \normalsize \textit{July'18-Present}  \\ 
Mentored 1$^{st} $-year students under the Student Mentor Programme conducted by Student Welfare Group.
\end{itemize}

\hspace{0.5cm}\\[-0.2cm]
\resheading{\textbf{AWARDS AND ACHIEVEMENTS} }\\[\lsep]
\\[-1em]
\begin{itemize}
\item Cleared regionals and participated in the nationals of Indian National Astronomy Olympiad conducted by HBCSE.
\\[\psep]
\item Secured state rank 2 in International Olympiad of English Language’ 15
conducted by SilverZone.
\\[\psep]
\item Recipient of the Goralal Syngal Memorial Scholarship.
\\[\psep]
\item Div. 1 rated Codechef user.
\end{itemize}

\end{document}
